% Experiment 18: Page Replacement Simulator

\section{Experiment 18: Page Replacement Simulator}

\subsection{Problem Statement}
The objective is to implement \texttt{pagerepl} to simulate FIFO, LRU, and OPT replacement.

The tool must:
\begin{itemize}
    \item Accept \verb|--frames <F>|.
    \item Read reference string length and page numbers from stdin.
    \item Simulate the three algorithms independently.
    \item Output fault counts and final frame contents.
\end{itemize}

\subsection{Source Code}
\begin{lstlisting}[language=C, caption={pagerepl.c}]
// Page Replacement Simulator
//pagerepl

#include <stdio.h>
#include <stdlib.h>
#include <string.h>

#define E_INPUT "E_INPUT"
#define E_RANGE "E_RANGE"

void print_error(const char *code, const char *msg) {
    fprintf(stderr, "ERROR: %s: %s\n", code, msg);
    exit(1);
}

int lookup(int *frames, int F, int page) {
    for (int i = 0; i < F; i++) {
        if (frames[i] == page) return i;
    }
    return -1;
}

void print_final(int *frames, int F) {
    printf("OK: FINAL");
    for (int i = 0; i < F; i++) printf(" %d", frames[i]);
    printf("\n");
}

void solve_fifo(int *pages, int L, int F) {
    printf("ALG FIFO\n");
    int *frames = malloc(sizeof(int) * F);
    for(int i=0; i<F; i++) frames[i] = -1;
    
    int faults = 0;
    int head = 0;

    for (int i = 0; i < L; i++) {
        int page = pages[i];
        if (lookup(frames, F, page) == -1) {
            faults++;
            frames[head] = page;
            head = (head + 1) % F;
        }
    }
    printf("OK: FAULTS %d\n", faults);
    print_final(frames, F);
    free(frames);
}

void solve_lru(int *pages, int L, int F) {
    printf("ALG LRU\n");
    int *frames = malloc(sizeof(int) * F);
    int *last_used = malloc(sizeof(int) * F); 
    for(int i=0; i<F; i++) { frames[i] = -1; last_used[i] = -1; }

    int faults = 0;
    int current_time = 0;

    for (int i = 0; i < L; i++) {
        int page = pages[i];
        int idx = lookup(frames, F, page);
        current_time++;

        if (idx != -1) {
            last_used[idx] = current_time;
        } else {
            faults++;
            int victim = -1;
            
            for (int j = 0; j < F; j++) {
                if (frames[j] == -1) {
                    victim = j;
                    break;
                }
            }

            if (victim == -1) {
                int min_time = 2000000000;
                for (int j = 0; j < F; j++) {
                    if (last_used[j] < min_time) {
                        min_time = last_used[j];
                        victim = j;
                    }
              
                }
            }
            
            frames[victim] = page;
            last_used[victim] = current_time;
        }
    }
    printf("OK: FAULTS %d\n", faults);
    print_final(frames, F);
    free(frames);
    free(last_used);
}

void solve_opt(int *pages, int L, int F) {
    printf("ALG OPT\n");
    int *frames = malloc(sizeof(int) * F);
    for(int i=0; i<F; i++) frames[i] = -1;

    int faults = 0;

    for (int i = 0; i < L; i++) {
        int page = pages[i];
        if (lookup(frames, F, page) == -1) {
            faults++;
            int victim = -1;

            for (int j = 0; j < F; j++) {
                if (frames[j] == -1) {
                    victim = j;
                    break;
                }
            }

            if (victim == -1) {
                int max_dist = -1;
                
                for (int j = 0; j < F; j++) {
                    int dist = 2000000000;
                    for (int k = i + 1; k < L; k++) {
                        if (pages[k] == frames[j]) {
                            dist = k;
                            break;
                        }
                    }
                    
                    if (dist > max_dist) {
                        max_dist = dist;
                        victim = j;
                    }
                  
                }
            }
            frames[victim] = page;
        }
    }
    printf("OK: FAULTS %d\n", faults);
    print_final(frames, F);
    free(frames);
}

int main(int argc, char *argv[]) {
    int F = 0;
    if (argc > 2 && strcmp(argv[1], "--frames") == 0) F = atoi(argv[2]);
    if (F < 1 || F > 64) print_error(E_RANGE, "frames must be 1..64");

    int L;
    if (scanf("%d", &L) != 1) return 0; 
    if (L < 0) print_error(E_INPUT, "bad length");

    int *pages = malloc(sizeof(int) * L);
    for (int i = 0; i < L; i++) {
        if (scanf("%d", &pages[i]) != 1) print_error(E_INPUT, "bad page");
        if (pages[i] < 0) print_error(E_RANGE, "page numbers must be >= 0");
    }

    solve_fifo(pages, L, F);
    solve_lru(pages, L, F);
    solve_opt(pages, L, F);

    free(pages);
    return 0;
}
\end{lstlisting}

\subsection{Sample Input and Output}

\textbf{Case 1:}
\begin{verbatim}
$ printf "12\n1 2 3 2 4 1 2 5 2 1 2 3" | ./pagerepl --frames 3
ALG FIFO
OK: FAULTS 8
OK: FINAL 5 3 2
ALG LRU
OK: FAULTS 7
OK: FINAL 3 2 1
ALG OPT
OK: FAULTS 6
OK: FINAL 3 2 5
\end{verbatim}

\textbf{Case 2:}
\begin{verbatim}
$ printf '2\n1 -1' | ./pagerepl --frames 2
ERROR: E_RANGE: page numbers must be >= 0
\end{verbatim}